% =====================================================================
%  2bat-ud1-11.tex  —  SESSIÓ 11: Full de càlcul: primer contacte
%  (sense preàmbul: s'inclou des de main.tex)
%  Criteris: C1 i C2. Sessió TECNOLÒGICA (guia pràctica de full de càlcul).
% =====================================================================
\horatitol{Sessió 11 — Full de càlcul: primer contacte amb matrices}%
{Reproduir amb el full de càlcul el que ja saps fer a mà: introduir matrices, sumar-les i aplicar un escalar.}

\subsection{Idea clau}

El full de càlcul (Excel o Google Sheets) és una eina ideal per treballar amb
matrices: una matriu \emph{és} una taula de cel·les organitzada en files i
columnes. En aquesta sessió no aprendràs res de matemàtiques nou; el que faràs és
\textbf{reproduir amb l'ordinador} operacions que ja domines a mà (sumar matrices
i multiplicar per un escalar) i \textbf{comprovar} que la màquina dona el mateix
resultat. Així guanyaràs confiança en l'eina i estalviaràs feina repetitiva.

\begin{idea}
\textbf{Vocabulari del full de càlcul}
\begin{itemize}[leftmargin=1.4em, itemsep=2pt]
  \item \textbf{Cel·la:} cada caseta, identificada per una lletra (columna) i un
        número (fila): \texttt{A1}, \texttt{B2}, \texttt{C3}\dots
  \item \textbf{Rang:} un bloc rectangular de cel·les, s'escriu amb dos punts:
        \texttt{A1:B2} és el bloc de $2\times 2$ que va de \texttt{A1} a \texttt{B2}.
  \item \textbf{Fórmula:} comença sempre amb el signe \texttt{=}. Per exemple,
        \texttt{=A1+D1} suma el contingut de dues cel·les.
  \item \textbf{Arrossegar:} després d'escriure una fórmula en una cel·la, es pot
        \emph{arrossegar} (des de la cantonada inferior dreta) per copiar-la a tot
        un bloc; les referències s'ajusten soles.
\end{itemize}
\end{idea}

\subsection{Exemples}

\begin{exemple}[on van les dades]
Per representar la matriu $A=\begin{pmatrix} 40 & 22 \\ 18 & 30 \end{pmatrix}$
en el full de càlcul, posem cada número en una cel·la:
\begin{center}
\begin{tabular}{c|cc}
 & \textbf{A} & \textbf{B} \\
\hline
\textbf{1} & 40 & 22 \\
\textbf{2} & 18 & 30 \\
\end{tabular}
\end{center}
La matriu ocupa el rang \texttt{A1:B2}. La cel·la \texttt{A1} és l'element de la
fila 1, columna 1 ($40$); la cel·la \texttt{B2}, el de la fila 2, columna 2 ($30$).
\end{exemple}

\begin{exemple}[la fórmula de la suma]
Si tenim la matriu $A$ al rang \texttt{A1:B2} i la matriu $B$ al rang
\texttt{D1:E2}, per sumar-les escrivim a la cel·la \texttt{A4} la fórmula
\[
\texttt{=A1+D1}
\]
i l'arrosseguem a tot el bloc \texttt{A4:B5}. Cada cel·la del resultat suma les
dues cel·les de la mateixa posició, exactament com la suma casella a casella que
fem a mà.
\end{exemple}

% ---------------------------------------------------------------------
\subsection{Guia pràctica (pas a pas)}

Treballaràs amb dues matrices que ja vas operar a mà:
\[
A=\begin{pmatrix} 40 & 22 \\ 18 & 30 \end{pmatrix},\qquad
B=\begin{pmatrix} 45 & 20 \\ 16 & 28 \end{pmatrix}.
\]

\begin{idea}
\textbf{Passos a l'ordinador}
\begin{enumerate}[leftmargin=1.8em, itemsep=3pt]
  \item Obre un full de càlcul nou. Escriu la matriu $A$ al rang \texttt{A1:B2}
        (un número a cada cel·la) i la matriu $B$ al rang \texttt{D1:E2}.
  \item \textbf{Suma.} A la cel·la \texttt{A4} escriu \texttt{=A1+D1} i prem Retorn.
        Després selecciona \texttt{A4} i \emph{arrossega} la fórmula fins a omplir
        el bloc \texttt{A4:B5}. Ja tens la matriu $A+B$.
  \item \textbf{Producte per un escalar.} Suposa que s'aplica una retallada del
        $10\%$ a la matriu suma (multiplicar per $0{,}9$). A la cel·la \texttt{A7}
        escriu \texttt{=A4*0,9} i arrossega-la al bloc \texttt{A7:B8}. Ja tens
        $0{,}9\cdot(A+B)$.
  \item \textbf{Comprova.} Calcula $A+B$ i $0{,}9\cdot(A+B)$ \emph{a mà} a la
        llibreta i compara-ho amb el que mostra la pantalla. Han de coincidir.
\end{enumerate}
\end{idea}

\textbf{Tasques per lliurar}
\begin{exercicis}
\begin{exlist}
  \item Reprodueix els passos anteriors amb les matrices $A$ i $B$ donades i
        anota a la llibreta la matriu $A+B$ i la matriu $0{,}9\cdot(A+B)$ que
        mostra el full de càlcul.
  \item Canvia un sol número de la matriu $A$ (per exemple, posa $50$ a la cel·la
        \texttt{A1}) i observa què passa amb els blocs de la suma i de l'escalar.
        Explica, en una frase, per què canvien sols.
\end{exlist}
\end{exercicis}

% ---------------------------------------------------------------------
\fullsolucions{Sessió 11}
\begin{solucions}
\begin{sollist}
  \item Resultats esperats al full de càlcul:
        \[ A+B=\begin{pmatrix} 85 & 42 \\ 34 & 58 \end{pmatrix},\qquad
           0{,}9\cdot(A+B)=\begin{pmatrix} 76{,}5 & 37{,}8 \\ 30{,}6 & 52{,}2 \end{pmatrix}. \]
        Han de coincidir amb el càlcul manual.
  \item En posar $50$ a \texttt{A1}, la cel·la \texttt{A4} (suma) passa a
        $50+45=95$ i la cel·la \texttt{A7} (escalar) a $0{,}9\per 95=85{,}5$.
        Canvien soles perquè les fórmules fan \emph{referència} a les cel·les
        d'origen: si canvia una dada, tot el que en depèn es recalcula
        automàticament.
\end{sollist}

\noindent\textbf{Nota per al professorat.} S'avalua la \emph{correcció tècnica}
(introduir les matrices en rangs, escriure i arrossegar les fórmules, usar
referències) i la \emph{interpretació} (entendre que el full de càlcul reprodueix
el càlcul manual i que el recàlcul automàtic ve de les referències). Aquesta
sessió alimenta l'instrument «Treballs pràctics i tecnològics» (15\%).
\end{solucions}
